% !TeX program = pdflatex
\documentclass[12pt,a4paper,oneside]{article}
\usepackage[T1]{fontenc}
\usepackage[utf8]{inputenc}
\usepackage{newtxtext,newtxmath}
\usepackage[spanish,es-nodecimaldot]{babel}
\usepackage[a4paper,top=2.25cm,bottom=2.2cm,left=2.4cm,right=2.2cm,headheight=28pt]{geometry}
\usepackage{microtype}
\usepackage{graphicx}
\usepackage{float}
\usepackage{booktabs,tabularx,array,longtable}
\usepackage{enumitem}
\usepackage{xcolor}
\usepackage{tikz}
\usetikzlibrary{calc,positioning}
\usepackage[most]{tcolorbox}
\usepackage{listings}
\usepackage{fancyhdr}
\usepackage{titlesec}
\usepackage{titletoc}
\usepackage{caption}
\usepackage{subcaption}
\usepackage{hyperref}
\usepackage{setspace}
\setstretch{1.25}
\setlength{\parindent}{0pt}
\setlength{\parskip}{7pt}
\definecolor{azul}{HTML}{122A46}
\definecolor{azulmedio}{HTML}{254B71}
\definecolor{turquesa}{HTML}{008D95}
\definecolor{dorado}{HTML}{E2B45E}
\definecolor{tint}{HTML}{F3F7FA}
\definecolor{gris}{HTML}{566579}
\definecolor{borde}{HTML}{D8E2EB}
\definecolor{codigo}{HTML}{F7F9FC}
% =================== DATOS EDITABLES ===================
\newcommand{\universidad}{UNIVERSIDAD TÉCNICA ESTATAL DE QUEVEDO}
\newcommand{\facultad}{FACULTAD DE CIENCIAS DE LA COMPUTACIÓN Y DISEÑOS DIGITALES}
\newcommand{\carrera}{CARRERA DE INGENIERÍA DE SOFTWARE}
\newcommand{\materia}{Administración de Base de Datos}
\newcommand{\tema}{Informe final: Repositorio de Tutorías}
\newcommand{\estudiante}{Kamila Calle, Roselyn Sánchez y Mayummy Trujillo}
\newcommand{\docente}{Cordero Bazurto Jose Steven}
\newcommand{\curso}{Quinto nivel — Paralelo A}
\newcommand{\fechaentrega}{22 de noviembre de 2026}
\newcommand{\lugar}{Quevedo -- Ecuador}
\newcommand{\anio}{2026}
\newcommand{\logoinstitucional}{\includegraphics[height=3.1cm,keepaspectratio]{Imagenes/EVIDENCIAS/Logo_UTEQ.png}}
% Comando para nombres de archivos de evidencia (usado en tablas)
\providecommand{\evus}{\textunderscore\allowbreak}
\providecommand{\ev}[1]{{\texttt{\scriptsize\let\_\evus #1}}}
\renewcommand{\thesection}{\arabic{section}.}
\renewcommand{\thesubsection}{\arabic{section}.\arabic{subsection}.}
\renewcommand{\thesubsubsection}{\arabic{section}.\arabic{subsection}.\arabic{subsubsection}.}
\hypersetup{unicode=true,colorlinks=true,linkcolor=black,urlcolor=black,citecolor=black,pdftitle={\tema},pdfauthor={\estudiante}}
\addto\captionsspanish{\renewcommand{\contentsname}{Contenido}\renewcommand{\listfigurename}{Índice de figuras}\renewcommand{\listtablename}{Índice de tablas}}
\titleformat{\section}{\Large\bfseries\color{black}}{\thesection}{0.8em}{}
\titleformat{\subsection}{\large\bfseries\color{black}}{\thesubsection}{0.75em}{}
\titleformat{\subsubsection}{\normalsize\bfseries\color{black}}{\thesubsubsection}{0.7em}{}
\titlespacing*{\section}{0pt}{19pt}{10pt}
\titlespacing*{\subsection}{0pt}{13pt}{5pt}
\titlecontents{section}[0em]{\vspace{4pt}\bfseries\color{black}}{\contentslabel{2.2em}}{}{\titlerule*[0.6pc]{.}\contentspage}
\captionsetup{font=small,labelfont={bf,color=black},labelsep=period}
\pagestyle{fancy}\fancyhf{}
\fancyhead[L]{\small\textbf{UTEQ}\enspace|\enspace Ingeniería de Software}
\fancyhead[R]{\small Informe académico}
\renewcommand{\headrulewidth}{0.6pt}
\renewcommand{\headrule}{\hbox to\headwidth{\color{borde}\leaders\hrule height \headrulewidth\hfill}}
\fancyfoot[C]{\small\thepage}
\renewcommand{\footrulewidth}{0pt}
\newtcolorbox{destacado}[1][]{enhanced,breakable,colback=tint,colframe=borde,boxrule=.6pt,arc=2mm,
  borderline west={3pt}{0pt}{turquesa},left=13pt,right=11pt,top=9pt,bottom=9pt,#1}
\newtcolorbox{resultado}[1][]{enhanced,breakable,colback=white,colframe=turquesa!55,boxrule=.75pt,arc=2mm,
 title={\bfseries Resultado o hallazgo},coltitle=black,colbacktitle=tint,#1}
\lstdefinestyle{estiloUTEQ}{backgroundcolor=\color{codigo},basicstyle=\ttfamily\footnotesize,
  keywordstyle=\color{black}\bfseries,commentstyle=\color{black}\itshape,
  stringstyle=\color{black},numbers=left,numberstyle=\tiny\color{black},
  numbersep=10pt,stepnumber=1,frame=single,rulecolor=\color{borde},framesep=8pt,
  breaklines=true,breakatwhitespace=false,showstringspaces=false,tabsize=2,
  keepspaces=true,columns=fullflexible,captionpos=b,xleftmargin=7pt,xrightmargin=3pt}
\lstset{style=estiloUTEQ}
\renewcommand{\lstlistingname}{Código}
\renewcommand{\lstlistlistingname}{Índice de códigos}
\begin{document}
% ========================= PORTADA =========================
\begin{titlepage}
\begin{tikzpicture}[remember picture,overlay]
  \fill[azul] (current page.north west) rectangle ([yshift=-1.05cm]current page.north east);
  \fill[turquesa] ([yshift=-1.05cm]current page.north west) rectangle ([yshift=-1.12cm]current page.north east);
  \fill[azul] ([yshift=1.05cm]current page.south west) rectangle (current page.south east);
  \fill[dorado] ([xshift=1.8cm,yshift=1.05cm]current page.south west) rectangle
  ([xshift=5.7cm,yshift=1.12cm]current page.south west);
\end{tikzpicture}
\vspace*{0.35cm}
\begin{center}
\logoinstitucional\\[0.55cm]
{\normalsize\bfseries\color{black}\universidad}\\[0.16cm]
{\small\color{black}\facultad}\\[0.12cm]
{\small\bfseries\color{black}\carrera}
\end{center}
\vspace{1.25cm}
\begin{center}
{\small\bfseries\color{black}\MakeUppercase{INFORME ACADÉMICO}}\\[0.48cm]
{\fontsize{23}{27}\selectfont\bfseries\color{black}\tema}\\[0.65cm]
{\color{dorado}\rule{3.7cm}{2pt}}\\[0.35cm]
{\large\color{black}\materia}
\end{center}
\vfill
\begin{tcolorbox}[colback=tint,colframe=borde,arc=2mm,boxrule=.7pt,left=15pt,right=15pt,top=12pt,bottom=12pt]
\renewcommand{\arraystretch}{1.75}
\begin{tabular}{@{}>{\bfseries\raggedright\arraybackslash}m{3.1cm}>{\raggedright\arraybackslash}m{10.7cm}@{}}
Estudiantes: & \estudiante\\
Docente: & \docente\\
Curso: & \curso\\
Entrega: & \fechaentrega
\end{tabular}
\end{tcolorbox}
\vspace{0.65cm}
\begin{center}{\bfseries\color{black}\lugar\quad |\quad\anio}\end{center}
\vspace{0.35cm}
\end{titlepage}
% ========================= PRELIMINARES =========================
\pagenumbering{roman}
\setcounter{page}{1}
\tableofcontents
\clearpage
\pagenumbering{arabic}
% ========================= CUERPO =========================
\section{Introducción}
Este informe documenta el desarrollo de la aplicación web \textit{Repositorio de Tutorías}, construida con C\# (ASP.NET Core MVC), Npgsql y PostgreSQL. El trabajo se dividió en tres módulos: RF-01 (periodos y categorías documentales), RF-02 (facultades, carreras y usuarios) y RF-03 (documentos). Todos los accesos a la base de datos usan consultas parametrizadas y las claves foráneas se definieron con \texttt{ON DELETE RESTRICT}.
\section{Misión}

El \textit{Repositorio de Tutorías} existe para guardar en un solo lugar, y clasificados
por periodo académico y por categoría documental, los documentos de las tutorías: el plan
de acción inicial, el informe de corte y el informe final.

Hoy esa información está repartida en correos electrónicos, archivos sueltos y
dispositivos personales. Cuando a un coordinador le piden un histórico de cinco años,
tiene que revisar cinco años de correos, y al no existir un lugar centralizado se pierde
información y cuesta encontrarla. Con la plataforma basta con cambiar de periodo y
consultar lo que se subió en ese semestre, en cualquier momento y desde cualquier lugar.

La plataforma no está pensada solo para la carrera de Software. Las facultades y las
carreras se registran como datos del sistema, de modo que cualquier facultad y cualquier
carrera de la institución pueda usarla con su propia información.

\section{Stakeholders}

La Tabla~\ref{tab:stakeholders} indica quiénes tienen interés en la plataforma o se ven
afectados por ella, y cuál es su relación con el sistema. Los tutores son quienes suben
los documentos; los demás roles organizan, consultan o evalúan.

\begin{longtable}{>{\raggedright\arraybackslash}p{3.7cm}>{\raggedright\arraybackslash}p{11.3cm}}
\caption{Stakeholders del sistema y su relación con él}\label{tab:stakeholders}\\
\toprule
\textbf{Stakeholder} & \textbf{Relación con el sistema} \\
\midrule
\endfirsthead
\toprule
\textbf{Stakeholder} & \textbf{Relación con el sistema} \\
\midrule
\endhead
Administrador general &
Configura la plataforma: registra facultades, carreras y usuarios, crea los periodos
y define qué categorías documentales se piden en cada semestre (por ejemplo, plan de
acción inicial, informe de corte e informe final). \\
\addlinespace
Coordinador de carrera &
Da seguimiento a los documentos de su carrera. Consulta lo que entregaron los tutores,
filtrando por periodo y categoría. \\
\addlinespace
Coordinador de área &
Apoya las operaciones de la aplicación junto con el administrador general y consulta
los documentos de las carreras que le corresponden. \\
\addlinespace
Tutores &
Son quienes suben los documentos. Eligen el periodo y la categoría, escriben el título
y adjuntan el archivo, que queda guardado en PostgreSQL. \\
\addlinespace
Docente de la materia (evaluador) &
Orienta y evalúa el proyecto. Revisa que la aplicación funcione, que la integración con
PostgreSQL sea correcta y que las decisiones estén documentadas. \\
\addlinespace
Institución (UTEQ) &
Es la beneficiaria final: conserva el histórico de tutorías ordenado y accesible,
sin depender de correos personales. \\
\bottomrule
\end{longtable}

\section{Objetivos}

\subsection{Objetivo general}

Desarrollar una aplicación web conectada a PostgreSQL que guarde y clasifique por
periodo académico y categoría documental los documentos de tutorías, para que pueda
ser usada por cualquier facultad y carrera y reemplace la búsqueda de documentos en
años de correos.

\subsection{Objetivos específicos}

\begin{enumerate}[label=\textbf{OE\arabic*.}, leftmargin=1.6cm]
  \item Gestionar periodos y categorías documentales, de modo que al crear un periodo
        nuevo se puedan reutilizar las categorías de periodos anteriores y agregar
        nuevas o quitar las que ya no se necesiten (módulo de periodos y categorías,
        RF-01).
  \item Registrar facultades, carreras y usuarios, asignando a cada usuario uno de los
        cuatro roles (administrador general, coordinador de carrera, coordinador de área
        o tutor), con validaciones de integridad como nombre único, correo válido y
        único, y carrera obligatoria para todos los roles excepto el administrador
        general (módulo de facultades, carreras y usuarios, RF-02).
  \item Permitir que los tutores suban documentos y que se consulten con filtros por
        periodo, categoría y carrera, con paginación, validando el tipo de archivo, el
        tamaño y que la categoría esté habilitada para ese periodo (módulo de
        documentos, RF-03).
  \item Diseñar la base de datos para que soporte millones de registros, usando índices
        en las llaves foráneas y en los campos de filtro, paginación en las listas y un
        script de datos de prueba masivos, para que las pantallas no se vuelvan lentas
        al crecer la información.
  \item Proteger la integridad de los datos con llaves foráneas \texttt{ON DELETE
        RESTRICT}, consultas parametrizadas y mensajes claros de éxito o de error en
        cada operación.
  \item Comprobar cada requisito con pruebas funcionales que registren los datos
        usados, el resultado esperado, el resultado obtenido y su evidencia, y que
        verifiquen en PostgreSQL lo que se hizo desde la aplicación.
\end{enumerate}

\section{Alcance}

\subsection{Qué incluye este semestre}

La primera versión de la aplicación es una aplicación web en C\# con ASP.NET Core MVC
y PostgreSQL como gestor de base de datos, donde se almacena toda la información. Incluye
tres módulos de prioridad alta (Tabla~\ref{tab:modulos}):

\begin{longtable}{>{\raggedright\arraybackslash}p{3.7cm}>{\raggedright\arraybackslash}p{11.3cm}}
\caption{Módulos incluidos en la primera versión}\label{tab:modulos}\\
\toprule
\textbf{Módulo} & \textbf{Funcionalidades incluidas} \\
\midrule
\endfirsthead
\toprule
\textbf{Módulo} & \textbf{Funcionalidades incluidas} \\
\midrule
\endhead
Periodos y categorías documentales (RF-01) &
Crear, listar, editar, cerrar y eliminar periodos; crear, editar y eliminar categorías
documentales; reutilizar categorías de periodos anteriores al crear un periodo nuevo.
Validaciones: nombre de periodo único, fecha de fin posterior a la de inicio, categoría
con nombre y no eliminar una categoría que ya tiene documentos. \\
\addlinespace
Facultades, carreras y usuarios (RF-02) &
Crear, listar, editar y eliminar facultades, carreras (cada una pertenece a una facultad)
y usuarios (con rol y carrera). Validaciones: nombre único, correo con formato válido y
único, carrera con facultad, carrera obligatoria salvo para el administrador general, y
no eliminar una facultad con carreras ni un usuario con documentos. \\
\addlinespace
Documentos (RF-03) &
Subir documentos en PDF o DOCX, consultarlos con filtros por periodo, categoría y carrera
con paginación y orden por fecha, editar sus datos, eliminarlos y descargarlos. \\
\bottomrule
\end{longtable}

Además, el alcance incluye la conexión de la aplicación con PostgreSQL, los scripts SQL
por requisito (\texttt{00\_base.sql}, \texttt{rf01.sql}, \texttt{rf02.sql} y
\texttt{rf03.sql}), las pruebas funcionales con evidencia y un script de datos de prueba
masivos para comprobar el rendimiento del diseño.

\subsection{Qué queda fuera}

\begin{itemize}[leftmargin=1.2cm]
  \item Inicio de sesión y control de acceso, que se contemplan para una versión posterior.
        Por ahora los roles se registran en el sistema, pero todavía no se restringen las
        pantallas según el rol de quien entra.
  \item Administración avanzada de roles y privilegios de PostgreSQL, que corresponde a
        contenidos posteriores de la materia.
  \item Carga real de millones de documentos. Solo se generan datos ficticios para
        probar que el diseño responde bien.
  \item Videollamadas, mensajería entre usuarios y agendamiento de tutorías, porque el
        problema que se resuelve es la dispersión y la pérdida de los documentos.
  \item Reportes y estadísticas de uso de la plataforma.
  \item Avisos o recordatorios por correo a los tutores.
  \item Importar automáticamente los documentos históricos que hoy están en correos.
\end{itemize}

\subsection{Cambios respecto a la propuesta inicial}

Después de la propuesta de la semana 2 el docente precisó cómo debe funcionar la
plataforma. Por eso se ajustó la propuesta inicial. La Tabla~\ref{tab:cambios} resume
los cambios y su motivo.

\begingroup\small
\begin{longtable}{>{\raggedright\arraybackslash}p{2.3cm}>{\raggedright\arraybackslash}p{3.6cm}>{\raggedright\arraybackslash}p{4.3cm}>{\raggedright\arraybackslash}p{3.9cm}}
\caption{Cambios respecto a la propuesta de la semana 2}\label{tab:cambios}\\
\toprule
\textbf{Elemento} & \textbf{Propuesta inicial} & \textbf{Versión actual} & \textbf{Motivo} \\
\midrule
\endfirsthead
\toprule
\textbf{Elemento} & \textbf{Propuesta inicial} & \textbf{Versión actual} & \textbf{Motivo} \\
\midrule
\endhead
Problema &
Materiales de tutorías para estudiantes y docentes. &
Documentos de tutorías (plan de acción inicial, informe de corte e informe final)
que hoy se buscan en años de correos. &
El docente explicó que el objetivo es evitar revisar cinco años de correos. \\
\addlinespace
Actores &
Administrador, docente y estudiante (tres actores). &
Administrador general, coordinador de carrera, coordinador de área y tutor (cuatro
roles). El estudiante ya no es un actor. &
El docente indicó estos roles; los tutores son quienes suben los documentos. \\
\addlinespace
Cobertura &
No se precisaba. &
Cualquier facultad y carrera, con facultades y carreras registradas como datos. &
El docente pidió que no sirva solo a la carrera de Software. \\
\addlinespace
Organización de los documentos &
Por asignatura, tema, tipo y fecha. &
Por periodo académico y categoría documental, reutilizable entre semestres. &
El administrador define cada semestre qué documentos se necesitan y las categorías
se reutilizan. \\
\addlinespace
Seguridad &
Inicio de sesión, control de acceso, cierre de sesión y contraseñas cifradas. &
Se deja para una versión posterior. Por ahora solo se registra el rol de cada usuario. &
Decisión del equipo para priorizar los tres módulos; la administración avanzada de
roles corresponde a contenidos posteriores. \\
\addlinespace
Rendimiento &
Búsquedas ágiles con al menos 1\,000 recursos. &
Diseño pensado para millones de registros, con índices, paginación y datos de prueba
masivos. &
El docente indicó que las bases de datos de la materia manejan millones de registros. \\
\bottomrule
\end{longtable}
\endgroup

\section{Desarrollo}
\subsection{Contexto y herramientas utilizadas}
Aplicación web en C\# con ASP.NET Core MVC y Npgsql, base de datos PostgreSQL ejecutada en Docker, y control de versiones con Git y GitHub (una rama por módulo y revisión cruzada de Pull Requests).

% ---------------------------- RF-01 ----------------------------
\subsection{RF-01: Periodos y categorías documentales}
Gestiona periodos (estado ACTIVO o CERRADO) y categorías documentales, y habilita categorías por periodo mediante la tabla \texttt{periodo\_categoria}.

\subsubsection{Dificultades técnicas y decisiones}
\begin{itemize}[leftmargin=1.8em,itemsep=3pt]
\item \textbf{Vista faltante.} Al abrir <<Nuevo periodo>> apareció el error de vista no encontrada, porque faltaba \texttt{Views/Periodos/Crear.cshtml}; se creó la vista.
\item \textbf{Transacción.} El periodo y sus categorías se guardan en una sola transacción, para no dejar periodos incompletos.
\item \textbf{Categorías sin duplicar.} El \texttt{INSERT} usa \texttt{ON CONFLICT}; si el nombre existe, se reutiliza la categoría.
\item \textbf{Claves foráneas.} Se usó \texttt{ON DELETE RESTRICT} y consultas parametrizadas en todo el módulo.
\item \textbf{Limitación.} La regla de no eliminar una categoría con documentos se probó primero solo contra \texttt{periodo\_categoria}; con RF-03 integrado quedó cubierta también por la tabla \texttt{documento}.
\end{itemize}

\subsubsection{Pruebas de RF-01}
\begin{table}[H]\centering
\caption{Pruebas funcionales de RF-01}\label{tab:rf01}
\footnotesize\renewcommand{\arraystretch}{1.3}
\begin{tabularx}{\linewidth}{@{}p{1.4cm}p{2.3cm}X X X p{2.4cm}@{}}
\toprule
Requisito & Prueba & Datos usados & Resultado esperado & Resultado obtenido & Evidencia\\ \midrule
RF-01 & Crear categoría documental & Nombre y descripción de la categoría; activa & Se guarda y aparece en la lista con mensaje de éxito & Categoría creada y visible en la lista & \ev{RF01\_categoria\_crear.png}\\
\addlinespace
RF-01 & Crear periodo con categorías & Nombre, fechas, estado ACTIVO y una categoría habilitada & Se guarda el periodo con sus categorías en una sola transacción & Periodo creado con su categoría habilitada & \ev{RF01\_periodo\_crear\_con\_categoria\_nueva.png}\newline\ev{RF01\_verificacion\_postgres\_periodo\_categoria.png}\\
\addlinespace
RF-01 & Categoría ya existente & Nombre de una categoría ya registrada & Se reutiliza la categoría, sin error ni duplicado & Categoría reutilizada mediante \texttt{ON CONFLICT} & \ev{RF01\_reutilizar\_categorias.png}\newline\ev{RF01\_reutilizar\_categorias\_postgres.png}\\
\addlinespace
RF-01 & Eliminar categoría en uso & Categoría habilitada en un periodo & No se elimina y se muestra un mensaje claro & Mensaje de que la categoría está en uso & \ev{RF01\_categoria\_eliminar.png}\\
\bottomrule
\end{tabularx}
\end{table}

% ---------------------------- RF-02 ----------------------------
\subsection{RF-02: Facultades, carreras y usuarios}
Gestiona facultades, carreras y usuarios, con roles y la carrera asociada a cada usuario.

\subsubsection{Dificultades técnicas y decisiones}
\begin{itemize}[leftmargin=1.8em,itemsep=3pt]
\item \textbf{Entorno incompleto.} La computadora no tenía Git ni el SDK de .NET 11; se instalaron antes de clonar el repositorio y compilar. También se recuperó la clave de PostgreSQL de la configuración del contenedor de Docker y se guardó con \texttt{dotnet user-secrets}, para que no quede escrita en ningún archivo del repositorio.
\item \textbf{Regla del administrador general.} La columna \texttt{usuario.id\_carrera} acepta vacío (NULL) porque el administrador general no pertenece a una sola carrera. Como la base de datos no puede distinguir por rol si el vacío es válido o no, se validó en el controlador de C\#: si el rol no es administrador general, la carrera es obligatoria. Con la misma regla, el coordinador de área también exige carrera.
\item \textbf{Nombre de carrera único por facultad.} El reparto solo pedía nombre único para la facultad. Para las carreras se usó una restricción \texttt{UNIQUE (id\_facultad, nombre)}: dos facultades pueden tener una carrera con el mismo nombre, pero una misma facultad no la repite.
\item \textbf{Integración con RF-01 y RF-03.} Al traer \texttt{main} a la rama no hubo conflictos y el proyecto compiló. Al revisar la rama de RF-03 se vio que su consulta de tutores usaba la columna \texttt{nombre}, pero en \texttt{usuario} se llama \texttt{nombres}, y que sus vistas estaban en \texttt{app/Documentos/} en vez de \texttt{app/Views/Documentos/}. Se comunicó para su corrección.
\item \textbf{Capturas.} DBeaver cerró la conexión por inactividad y mostró <<Not connected to database>>; se reconectó y se repitió la captura.
\end{itemize}

\subsubsection{Pruebas de RF-02}
Cada prueba se hizo desde la aplicación y después se comprobó en PostgreSQL con una consulta, para confirmar que lo que se vio en pantalla coincide con lo guardado en la base de datos. Las capturas están en la carpeta \texttt{evidencias/} del repositorio.

\begingroup\footnotesize
\setlength{\tabcolsep}{3pt}
\begin{longtable}{>{\raggedright\arraybackslash}p{1.5cm}>{\raggedright\arraybackslash}p{1.8cm}>{\raggedright\arraybackslash}p{2.8cm}>{\raggedright\arraybackslash}p{2.7cm}>{\raggedright\arraybackslash}p{2.7cm}>{\raggedright\arraybackslash}p{3.2cm}}
\caption{Pruebas funcionales de RF-02}\label{tab:pruebas-rf02}\\
\toprule
\textbf{Requisito} & \textbf{Prueba} & \textbf{Datos usados} & \textbf{Resultado esperado} & \textbf{Resultado obtenido} & \textbf{Evidencia} \\
\midrule
\endfirsthead
\toprule
\textbf{Requisito} & \textbf{Prueba} & \textbf{Datos usados} & \textbf{Resultado esperado} & \textbf{Resultado obtenido} & \textbf{Evidencia} \\
\midrule
\endhead
RF-02 & Crear facultad &
Nombre: Facultad de Ciencias de la Ingeniería; activa. &
Se guarda y aparece en la lista con el mensaje ``Facultad creada''. &
Apareció el mensaje y la facultad en la lista. En PostgreSQL, \texttt{facultad} tiene la fila con \texttt{id\_facultad = 1}. &
\ev{RF02\_crear\_facultad.png}\newline
\ev{RF02\_crear\_facultad\_sql.png} \\
\addlinespace
RF-02 & Crear carrera &
Nombre: Ingeniería de Software; facultad elegida de la lista. &
Se guarda asociada a su facultad, con el mensaje ``Carrera creada''. &
Apareció el mensaje y la carrera junto a su facultad. En PostgreSQL, la consulta con \texttt{JOIN} muestra la carrera con su facultad. &
\ev{RF02\_crear\_carrera.png}\newline
\ev{RF02\_crear\_carrera\_sql.png} \\
\addlinespace
RF-02 & Crear usuario tutor &
Nombres: Tutor Prueba Uno; correo: tutor1@uteq.edu.ec; rol: tutor; carrera: Ingeniería de Software. &
Se guarda con su rol y carrera, con el mensaje ``Usuario creado''. &
Apareció el mensaje y el usuario en la lista. En PostgreSQL, \texttt{usuario} muestra el tutor con su rol y su carrera. &
\ev{RF02\_crear\_usuario\_tutor.png}\newline
\ev{RF02\_crear\_usuario\_tutor\_sql.png} \\
\addlinespace
RF-02 & Error: correo repetido &
Nombres: Tutor Prueba Uno; correo: tutor1@uteq.edu.ec (ya registrado); rol: tutor; misma carrera. &
No se guarda y se muestra ``Ya existe un usuario con ese correo''. &
El sistema rechazó el registro y mostró el mensaje bajo el campo Correo. En PostgreSQL, el correo aparece una sola vez. &
\ev{RF02\_error\_correo\_repetido.png}\newline
\ev{RF02\_error\_correo\_repetido\_sql.png} \\
\addlinespace
RF-02 & Error: eliminar facultad con carreras &
Facultad de Ciencias de la Ingeniería, que tiene 1 carrera asociada. &
No se elimina y se muestra ``No se puede eliminar: la facultad tiene carreras asociadas''. &
El sistema mostró el mensaje y la facultad siguió en la lista. En PostgreSQL, la facultad sigue con 1 carrera. &
\ev{RF02\_error\_eliminar\_facultad.png}\newline
\ev{RF02\_error\_eliminar\_facultad\_sql.png} \\
\addlinespace
RF-02 & Error: tutor sin carrera &
Nombres: Tutor Sin Carrera; correo: tutor3@uteq.edu.ec; rol: tutor; sin carrera. &
No se guarda y se muestra ``Solo el administrador general puede no tener carrera''. &
El sistema rechazó el registro y mostró el mensaje bajo el campo Carrera. &
\ev{RF02\_error\_tutor\_sin\_carrera.png} \\
\bottomrule
\end{longtable}
\endgroup

Las consultas usadas para comprobar los datos en PostgreSQL fueron:

\begingroup\sloppy\raggedright
\begin{itemize}[leftmargin=1.2cm]
  \item \ev{SELECT * FROM facultad;}
  \item \ev{SELECT c.id\_carrera, c.nombre, f.nombre AS facultad FROM carrera c JOIN
        facultad f ON f.id\_facultad = c.id\_facultad;}
  \item \ev{SELECT u.id\_usuario, u.nombres, u.correo, r.nombre AS rol, c.nombre AS
        carrera FROM usuario u JOIN rol r ON r.id\_rol = u.id\_rol LEFT JOIN carrera c ON
        c.id\_carrera = u.id\_carrera;}
  \item \ev{SELECT correo, COUNT(*) AS veces FROM usuario GROUP BY correo;}
  \item \ev{SELECT f.id\_facultad, f.nombre, COUNT(c.id\_carrera) AS carreras FROM
        facultad f LEFT JOIN carrera c ON c.id\_facultad = f.id\_facultad GROUP BY
        f.id\_facultad, f.nombre;}
\end{itemize}
\endgroup

Con el módulo de documentos (RF-03) ya integrado en la rama principal, se comprobó que no se puede eliminar un usuario, una categoría, una carrera ni un periodo que tengan documentos asociados.

% ---------------------------- RF-03 ----------------------------
\subsection{RF-03: Documentos}
Se creó la tabla \texttt{documento} con llaves foráneas hacia \texttt{usuario}, \texttt{periodo}, \texttt{categoria\_documental} y \texttt{carrera}, todas con \texttt{ON DELETE RESTRICT}. Se definieron índices sobre (\texttt{id\_periodo}, \texttt{id\_categoria}), \texttt{id\_carrera} y \texttt{fecha\_subida}. Las validaciones incluyen: solo PDF y DOCX, tamaño máximo de 10 MB, título obligatorio, categoría habilitada para el periodo y periodo cerrado sin documentos nuevos. La lista permite filtrar por periodo, categoría y carrera, con paginación (\texttt{LIMIT} y \texttt{OFFSET}) y orden por fecha.

\begin{lstlisting}[language=SQL,caption={Tabla documento (resumen)},label={lst:doc}]
CREATE TABLE IF NOT EXISTS documento (
    id_documento SERIAL PRIMARY KEY,
    titulo VARCHAR(200) NOT NULL,
    nombre_archivo VARCHAR(255) NOT NULL,
    ruta_archivo VARCHAR(500) NOT NULL,
    tipo_archivo VARCHAR(10) NOT NULL,
    tamano_bytes INTEGER NOT NULL,
    fecha_subida TIMESTAMP NOT NULL DEFAULT NOW(),
    id_usuario INTEGER NOT NULL REFERENCES usuario(id_usuario) ON DELETE RESTRICT,
    id_periodo INTEGER NOT NULL REFERENCES periodo(id_periodo) ON DELETE RESTRICT,
    id_categoria INTEGER NOT NULL REFERENCES categoria_documental(id_categoria) ON DELETE RESTRICT,
    id_carrera INTEGER NOT NULL REFERENCES carrera(id_carrera) ON DELETE RESTRICT
);
\end{lstlisting}

\subsubsection{Dificultades técnicas}
\begin{itemize}[leftmargin=1.8em,itemsep=3pt]
\item Las vistas estaban fuera de \texttt{Views} y la aplicación no las encontraba; se movieron a \texttt{Views/Documentos}.
\item La columna del tutor se llama \texttt{nombres} y no \texttt{nombre}; además se filtró el rol tutor.
\item La versión del SDK de .NET difería entre equipos (.NET 10 frente a .NET 11).
\item Los comandos de Git no funcionaban en la terminal de Visual Studio; se usaron los menús del entorno.
\end{itemize}

\subsubsection{Pruebas de RF-03}
\begin{table}[H]\centering
\caption{Pruebas de RF-03}\label{tab:rf03}
\small\renewcommand{\arraystretch}{1.3}
\begin{tabularx}{\linewidth}{@{}p{2.1cm}X X X p{2.6cm}@{}}
\toprule
Prueba & Datos usados & Esperado & Obtenido & Evidencia\\ \midrule
Subir documento & Título <<Plan de tutorias 2026>>, PDF, periodo 2026-2027 & Aparece <<Documento subido>> & Correcto & Figura~\ref{fig:subir}\\
Filtrar & Periodo 2026-2027 y categoría Planificación de tutorías & Solo documentos coincidentes & Correcto & Figura~\ref{fig:filtrar}\\
Editar & Nuevo título & Aparece <<Documento actualizado>> & Correcto & Figura~\ref{fig:editar}\\
Archivo no permitido & Archivo .txt & Rechazo con mensaje & Correcto & Figura~\ref{fig:error}\\
\bottomrule
\end{tabularx}
\end{table}

\begin{figure}[H]\centering
\includegraphics[width=.8\linewidth]{Imagenes/EVIDENCIAS/RF03_tablas_creadas.png}
\caption{Tablas creadas en PostgreSQL.}
\end{figure}

\begin{figure}[H]\centering
\includegraphics[width=.8\linewidth]{Imagenes/EVIDENCIAS/RF03_subir_documento.png}
\caption{Subida de documentos.}\label{fig:subir}
\end{figure}

\begin{figure}[H]\centering
\includegraphics[width=.8\linewidth]{Imagenes/EVIDENCIAS/RF03_subir_documento_sql.png}
\caption{Comprobación en PostgreSQL de que el documento se guardó (\texttt{SELECT * FROM documento;}).}
\end{figure}

\begin{figure}[H]\centering
\includegraphics[width=.8\linewidth]{Imagenes/EVIDENCIAS/RF03_filtrar_periodo_categoria.png}
\caption{Filtro por periodo y categoría.}\label{fig:filtrar}
\end{figure}

\begin{figure}[H]\centering
\includegraphics[width=.8\linewidth]{Imagenes/EVIDENCIAS/RF03_editar_documento.png}
\caption{Edición de un documento.}\label{fig:editar}
\end{figure}

\begin{figure}[H]\centering
\includegraphics[width=.6\linewidth]{Imagenes/EVIDENCIAS/RF03_error_archivo_no_permitido.png}
\caption{Rechazo de un archivo no permitido.}\label{fig:error}
\end{figure}

% ---------------------- RENDIMIENTO (NUEVO) ----------------------
\subsection{Rendimiento con datos masivos}
Para comprobar el diseño con más datos se cargaron 200{,}000 documentos de prueba y se midieron tres consultas con \texttt{EXPLAIN ANALYZE}, primero sin índices y después con índices. La tabla siguiente resume los tiempos y el plan de ejecución de cada consulta. Estas pruebas se hicieron con 200{,}000 documentos y no con millones de registros.

\begin{table}[H]\centering
\caption{Rendimiento con 200{,}000 documentos (\texttt{EXPLAIN ANALYZE})}\label{tab:rendimiento}
\small\renewcommand{\arraystretch}{1.3}
\begin{tabularx}{\linewidth}{@{}X r r r r X@{}}
\toprule
Consulta & Filas & Sin índices (ms) & Con índices (ms) & Mejora & Cambio de plan\\ \midrule
Conteo total & 1 & 32.175 & 27.064 & 1.2x & ---\\
\addlinespace
Filtro por carrera & 20{,}000 & 49.460 & 14.548 & 3.4x & \textit{Seq Scan} a \textit{Bitmap Index Scan}\\
\addlinespace
Filtro por periodo y categoría & 2{,}502 & 22.062 & 3.149 & 7.0x & \textit{Parallel Seq Scan} a \textit{Bitmap Index Scan}\\
\bottomrule
\end{tabularx}
\end{table}

Los índices mejoraron las consultas con filtro: el filtro por carrera pasó de un recorrido secuencial (\textit{Seq Scan}) a un \textit{Bitmap Index Scan}, y el filtro por periodo y categoría pasó de un \textit{Parallel Seq Scan} a un \textit{Bitmap Index Scan}. En el conteo total la mejora es pequeña, porque esa consulta debe contar todas las filas.

\begin{figure}[H]\centering
\includegraphics[width=.8\linewidth]{Imagenes/EVIDENCIAS/Rendimiento_lista_documentos.png}
\caption{Rendimiento: lista de documentos con los datos de prueba.}
\end{figure}

\begin{figure}[H]\centering
\includegraphics[width=.8\linewidth]{Imagenes/EVIDENCIAS/Rendimiento_sin_indices.png}
\caption{Rendimiento: consultas sin índices (\texttt{EXPLAIN ANALYZE}).}
\end{figure}

\begin{figure}[H]\centering
\includegraphics[width=.8\linewidth]{Imagenes/EVIDENCIAS/Rendimiento_con_indices.png}
\caption{Rendimiento: consultas con índices (\texttt{EXPLAIN ANALYZE}).}
\end{figure}

\subsection{Análisis de los resultados}
\begin{resultado}
Las cuatro pruebas de RF-03 finalizaron correctamente. Además, al integrar los tres módulos se comprobó que la base de datos impide eliminar usuarios, categorías, carreras o periodos que tengan documentos asociados, mostrando un mensaje claro.

Las pruebas de rendimiento se hicieron con 200{,}000 documentos y no con millones: con índices, el filtro por periodo y categoría fue 7.0 veces más rápido (de 22.062 ms a 3.149 ms) y el filtro por carrera, 3.4 veces más rápido.
\end{resultado}

\section{Limitaciones}
\begin{itemize}[leftmargin=1.8em,itemsep=3pt]
\item No hay inicio de sesión; queda para una versión posterior.
\item Solo se aceptan archivos PDF y DOCX de hasta 10 MB.
\item Los archivos se guardan en el servidor, sin copias en la nube.
\item Al editar un documento no se puede cambiar el archivo subido.
\item No hay versionado ni búsqueda por el contenido de los documentos.
\item Un periodo CERRADO no recibe documentos nuevos.
\item El rendimiento se probó con 200{,}000 documentos y no con millones de registros.
\item \textbf{Tiempo:} el proyecto se desarrolló en un semestre y con entregas semanales, por lo que se priorizaron los tres módulos principales y la seguridad (inicio de sesión) se dejó para una versión posterior.
\item \textbf{Costo:} se usaron herramientas gratuitas (PostgreSQL, Docker, GitHub y Overleaf), sin presupuesto para servidores, almacenamiento en la nube ni servicios de pago.
\item \textbf{Recursos:} las pruebas se hicieron en computadoras personales de las integrantes, no en un servidor de producción, y el equipo es de tres personas, con un módulo por integrante.
\end{itemize}

\section{Riesgos}
\begin{table}[H]\centering
\caption{Riesgos identificados}
\small\renewcommand{\arraystretch}{1.3}
\begin{tabularx}{\linewidth}{@{}p{0.6cm}X p{1.8cm} X@{}}
\toprule
\# & Riesgo & Impacto & Mitigación\\ \midrule
1 & Pérdida de archivos si falla el disco del servidor & Alto & Respaldos periódicos de archivos y base de datos\\
2 & Archivos dañinos con extensión válida & Alto & Validar tipo y tamaño; antivirus a futuro\\
3 & Acceso sin control por no existir inicio de sesión & Alto & Uso en red interna; autenticación en la siguiente versión\\
4 & Lentitud al crecer la cantidad de documentos & Medio & Índices y paginación\\
5 & Datos inconsistentes al eliminar registros relacionados & Alto & \texttt{ON DELETE RESTRICT} y mensajes claros\\
6 & Desajustes entre ramas al integrar el código & Medio & Convención de nombres y revisión cruzada de Pull Requests\\
\bottomrule
\end{tabularx}
\end{table}

\section{Roles de usuario}
\begin{itemize}[leftmargin=1.8em,itemsep=3pt]
\item \textbf{Administrador general:} administra facultades, carreras, usuarios, periodos y categorías documentales, y define qué documentos se piden en cada semestre. No pertenece a una sola carrera.
\item \textbf{Coordinador de carrera:} da seguimiento a los documentos y tutores de su carrera, consultando por periodo y categoría para verificar qué entregó cada tutor.
\item \textbf{Coordinador de área:} revisa los documentos de las carreras de su área y apoya al administrador general en las operaciones de la aplicación.
\item \textbf{Tutor:} sube los documentos de tutoría (plan de acción inicial, informe de corte e informe final), los edita y los consulta dentro de su carrera.
\end{itemize}

\section{Conclusiones}
\begin{enumerate}[leftmargin=1.8em]
\item Los tres módulos funcionan de forma integrada en la rama \texttt{main}.
\item Las restricciones \texttt{ON DELETE RESTRICT} evitan eliminar registros con documentos asociados.
\item Con 200{,}000 documentos de prueba, los índices mejoraron las consultas con filtro (hasta 7.0 veces); no se probó con millones de registros.
\item Como mejora futura se recomienda añadir inicio de sesión, respaldos automáticos y almacenamiento en la nube.
\end{enumerate}

\appendix
\section{Anexos}
Repositorio del proyecto: \url{https://github.com/mtrujillov-sys/Proyecto_De_AdministracionBD_Grupo07}. Las capturas de cada módulo están en la carpeta \texttt{evidencias}.
\end{document}